\begin{thebibliography}{25}
\bibitem[1]{ref1} Michael Alekhnovich, Eli Ben-Sasson, Alexander A. Razborov, and Avi Wigderson. Space complexity in propositional calculus. SIAM J. Comput., 31(4):1184–1211, 2002. doi:10.1137/S0097539700366735.
\bibitem[2]{ref2} Michael Alekhnovich, Eli Ben-Sasson, Alexander A. Razborov, and Avi Wigderson. Pseudor- andom generators in propositional proof complexity. SIAM J. Comput., 34(1):67–88, 2004. doi:10.1137/S0097539701389944.
\bibitem[3]{ref3} Michael Alekhnovich and Alexander A. Razborov. Lower bounds for polynomial calculus: Non- binomial case. Proceedings of the Steklov Institute of Mathematics, 242:18–35, 2003. Available at http://people.cs.uchicago.edu/~razborov/files/misha.pdf.Preliminary version in FOCS ’01.
\bibitem[4]{ref4} Albert Atserias and Víctor Dalmau. A combinatorial characterization of resolution width. J. Comput. Syst. Sci., 74(3):323–334, 2008. doi:10.1016/j.jcss.2007.06.025.
\bibitem[5]{ref5} Paul Beame and Toniann Pitassi. Propositional proof complexity: Past, present and future. Bull. EATCS, 65:66–89, 1998.
\bibitem[6]{ref6} Eli Ben-Sasson and Avi Wigderson. Short proofs are narrow - resolution made simple. J. ACM, 48(2):149–169, 2001. doi:10.1145/375827.375835.
\bibitem[7]{ref7} Matthew Clegg, Jeff Edmonds, and Russell Impagliazzo. Using the groebner basis algorithm to find proofs of unsatisfiability. In Proceedings of the Twenty-Eighth Annual ACM Symposium on the Theory of Computing, Philadelphia, Pennsylvania, USA, May 22-24, 1996, pages 174–183, 1996. doi:10.1145/237814.237860.
\bibitem[8]{ref8} Stephen A. Cook and Robert A. Reckhow. The relative efficiency of propositional proof systems. J. Symb. Log., 44(1):36–50, 1979. doi:10.2307/2273702.
\bibitem[9]{ref9} Susanna F. de Rezende, Jakob Nordström, Kilian Risse, and Dmitry Sokolov. Exponential resolution lower bounds for weak pigeonhole principle and perfect matching formulas over sparse graphs. In Shubhangi Saraf, editor, 35th Computational Complexity Conference, CCC 2020, July 28-31, 2020, Saarbrücken, Germany (Virtual Conference), volume 169 of LIPIcs, pages 28:1–28:24. Schloss Dagstuhl - Leibniz-Zentrum für Informatik, 2020. doi:10.4230/LIPIcs.CCC.2020.28.
\bibitem[10]{ref10} Konstantinos Georgiou, Avner Magen, and Madhur Tulsiani. Optimal sherali-adams gaps from pairwise independence. In Irit Dinur, Klaus Jansen, Joseph Naor, and José Rolim, editors, Approximation, Randomization, and Combinatorial Optimization. Algorithms and Techniques, pages 125–139, Berlin, Heidelberg, 2009. Springer Berlin Heidelberg.
\bibitem[11]{ref11} Russell Impagliazzo, Pavel Pudlák, and Jirí Sgall. Lower bounds for the polynomial calculus and the gröbner basis algorithm. Computational Complexity, 8(2):127–144, 1999. doi:10.1007/s000370050024.
\bibitem[12]{ref12} Jan Krajíček. Bounded arithmetic, propositional logic, and complexity theory, volume 60 of Encyclopedia of mathematics and its applications. Cambridge University Press, 1995.
\bibitem[13]{ref13} Jan Krajíček. On the weak pigeonhole principle. Fundamenta Mathematicae, 170:123–140, January 2001. doi:10.4064/fm170-1-8.
\bibitem[14]{ref14} Jan Krajíček. Dual weak pigeonhole principle, pseudo-surjective functions, and provability of circuit lower bounds. J. Symb. Log., 69(1):265–286, 2004. doi:10.2178/jsl/1080938841.
\bibitem[15]{ref15} Jan Krajíček. Proof complexity generators. In Hajo Broersma, Stefan S. Dantchev, Matthew Johnson, and Stefan Szeider, editors, Algorithms and Complexity in Durham 2006 - Proceedings of the Second ACiD Workshop, 18-20 September 2006, Durham, UK, volume 7 of Texts in Algorithmics, page 3. King’s College, London, 2006.
\bibitem[16]{ref16} Noam Nisan and Avi Wigderson. Hardness vs randomness. J. Comput. Syst. Sci., 49(2):149–167, 1994. doi:10.1016/S0022-0000(05)80043-1.
\bibitem[17]{ref17} Pavel Pudlák. Proofs as games. Am. Math. Mon., 107(6):541–550, 2000. URL: http: //www.jstor.org/stable/2589349.
\bibitem[18]{ref18} Alexander A. Razborov. Bounded arithmetic and lower bounds in boolean complexity. In Peter Clote and Jeffrey B. Remmel, editors, Feasible Mathematics II, pages 344–386, Boston, MA, 1995. Birkhäuser Boston.
\bibitem[19]{ref19} Alexander A. Razborov. Lower bounds for propositional proofs and independence results in bounded arithmetic. In Friedhelm Meyer auf der Heide and Burkhard Monien, editors, Automata, Languages and Programming, 23rd International Colloquium, ICALP96, Paderborn, Germany, 8-12 July 1996, Proceedings, volume 1099 of Lecture Notes in Computer Science, pages 48–62. Springer, 1996. doi:10.1007/3-540-61440-0\_116.
\bibitem[20]{ref20} Alexander A. Razborov. Improved resolution lower bounds for the weak pigeonhole prin- ciple. Electron. Colloquium Comput. Complex., 8(55), 2001. URL: http://eccc.hpi-web.de/eccc-reports/2001/TR01-055/index.html.
\bibitem[21]{ref21} Alexander A. Razborov. Resolution lower bounds for the weak functional pigeonhole principle. Theor. Comput. Sci., 303(1):233–243, 2003. doi:10.1016/S0304-3975(02)00453-X.
\bibitem[22]{ref22} Alexander A. Razborov. Pseudorandom generators hard for k-dnf resolution and polynomial calculus resolution. Ann. of Math., 181:415–472, 2015. doi:10.4007/annals.2015.181.2.1.
\bibitem[23]{ref23} Nathan Segerlind, Samuel R. Buss, and Russell Impagliazzo. A switching lemma for small restrictions and lower bounds for k-dnf resolution. SIAM J. Comput., 33(5):1171–1200, 2004. doi:10.1137/S0097539703428555.
\bibitem[24]{ref24} Dmitry Sokolov. (Semi)Algebraic proofs over $\pm$1 variables. In Konstantin Makarychev, Yury Makarychev, Madhur Tulsiani, Gautam Kamath, and Julia Chuzhoy, editors, Proccedings of the 52nd Annual ACM SIGACT Symposium on Theory of Computing, STOC 2020, Chicago, IL, USA, June 22-26, 2020, pages 78–90. ACM, 2020. doi:10.1145/3357713.3384288.
\bibitem[25]{ref25} Andrew Chi-Chih Yao. Theory and applications of trapdoor functions (extended abstract). In 23rd Annual Symposium on Foundations of Computer Science, Chicago, Illinois, USA, 3-5 November 1982, pages 80–91. IEEE Computer Society, 1982. doi:10.1109/SFCS.1982.45.
\end{thebibliography}
